\section{Requirement Engineering}

The system requirements are divided into functional requirements, which define the specific behaviors and services the platform must provide, and non-functional requirements, which specify quality attributes such as scalability, portability, and security.

\subsection{Functional Requirements}

The platform is designed to eliminate tool fragmentation by consolidating essential database operations into a single unified lifecycle.

\subsubsection*{Unified Management (Multi-Model Support)}
\begin{itemize}
    \item The system shall provide a centralized dashboard capable of managing both Relational (SQL) and Non-Relational (NoSQL) databases simultaneously.
    \item The platform must support PostgreSQL for structured relational data and MongoDB for flexible, document-based data.
    \item \textbf{Justification:} This requirement supports a best-of-breed selection approach, recognizing that different workloads require different database engines (e.g., SQL for transactional integrity and NoSQL for flexible schemas) rather than a monolithic solution.
\end{itemize}

\subsubsection*{Schema Operations \& Visual Design}
\begin{itemize}
    \item \textbf{Visual Builder:} The system shall include a drag-and-drop or interactive graphical interface for creating tables and defining entity relationships (ERD), abstracting complex syntax from the user.
    \item \textbf{File Import:} The system shall allow users to initialize projects by uploading existing CSV or JSON files, automatically inferring the initial schema structure.
    \item \textbf{LLM Agent-Powered Generation:} The platform shall implement an autonomous LLM agent module capable of converting natural language prompts (e.g., ``Create a schema for an e-commerce store'') into executable database structures, suggesting normalized relations, indexes, and validation rules, and producing visual ERD artifacts for user review.
    \item \textbf{Interactive Refinement:} The LLM agent shall support iterative, conversational refinement of generated schemas (user can accept, modify, or ask the agent to refactor), preserving audit logs of changes for reproducibility.
\end{itemize}

\subsubsection*{Data Manipulation \& Query Execution}
\begin{itemize}
    \item \textbf{Table Interface:} The system shall provide a spreadsheet-like interface for performing CRUD (Create, Read, Update, Delete) operations without requiring manual query writing.
    \item \textbf{Unified Query Editor:} The system shall support execution of both SQL queries and NoSQL aggregation pipelines within a single integrated environment.
    \item \textbf{Agent-Assisted Query Generation:} An LLM agent shall translate user natural language requests into syntactically correct SQL or aggregation pipelines, explain query semantics, and propose optimized alternatives (e.g., index hints, projection reduction) when applicable.
    \item \textbf{Context:} These features directly address ecosystem fragmentation by eliminating the need for developers to switch between isolated tools for data entry and query execution.
\end{itemize}

\subsubsection*{AI Assistance (LLM Agent-Centric)}
\begin{itemize}
    \item \textbf{Conversational Assistant:} The system shall integrate a persistent LLM agent that is schema-aware (has canonicalized knowledge of the current project schema and metadata) and can answer questions, generate queries, and propose schema edits grounded in the actual database contents and structure.
    \item \textbf{Explainability \& Traceability:} The LLM agent shall produce human-readable rationales for any generated schema change or query and shall attach links to the specific schema elements, sample rows, or tool outputs that informed its suggestion.
    \item \textbf{Multi-Agent Workflows:} For complex tasks (e.g., data migration, normalization, or cross-store joins), the platform shall support orchestrated LLM agents with specialized roles (planner, verifier, executor) to propose plans, validate transformations, and execute safe operations under policy constraints.
    \item \textbf{Human-in-the-Loop Controls:} The system shall allow users to require explicit approval before any destructive operation suggested by an agent (e.g., DROP, mass UPDATE), and record approvals in an auditable history.
    \item \textbf{Market Alignment:} This requirement aligns with the industry trend toward LLM-driven developer tooling while emphasizing agentic correctness, provenance, and human oversight rather than generic retrieval pipelines.
\end{itemize}

\subsubsection*{Cloud Storage \& Versioning}
\begin{itemize}
    \item \textbf{State Management:} The system shall securely store database projects in the cloud and maintain version histories that enable rollback to previous states.
    \item \textbf{Export \& Restore:} Users shall be able to export schemas and data or restore projects from saved versions to ensure data resilience and recovery.
    \item \textbf{Reliability:} Automated backup and restore mechanisms shall be implemented to reduce administrative overhead and improve service reliability.
    \item \textbf{Agent-Aware Checkpoints:} The platform shall capture agent proposals and transformations as part of version history (including agent prompts, outputs, and validation results) to enable reproducible rollbacks and post-hoc audits.
\end{itemize}
\subsubsection{Frontend Functional Requirements}
The frontend system is designed to provide an intuitive interface for managing DBaaS clusters. Key functional requirements include:
\begin{itemize}
    \item \textbf{Dynamic Schema Visualization:} Integration of a Mermaid-based rendering engine to dynamically generate and display Entity-Relationship Diagrams (ERDs) based on AI outputs.
    \item \textbf{Data Presentation Layer:} Support for server-side pagination, sorting, and filtering for tabular data display to ensure performance with large datasets.
    \item \textbf{Interactive Dashboard:} Implementation of real-time charts to visualize database health, resource consumption, and telemetry metrics.
    \item \textbf{Structured Data Handling:} specialized UI components for rendering JSON data in a clear, hierarchical, and readable format.
    \item \textbf{Project Management:} Tools for importing and exporting project configurations to ensure portability and data integrity.
\end{itemize}
\subsection{Non-Functional Requirements}

These requirements ensure that the system remains scalable, portable, secure, and suitable for production deployment.

\subsubsection*{Scalability via Containerization}
\begin{itemize}
    \item \textbf{Architecture:} The system shall utilize Docker for containerization.
    \item \textbf{Resource Efficiency:} Containers shall share the host operating system kernel, reducing resource overhead and enabling faster startup times compared to Virtual Machines (VMs).
    \item \textbf{Justification:} Although VMs provide stronger isolation for I/O-intensive workloads, containerization is selected for its superior efficiency in PaaS and SaaS environments, enabling dynamic horizontal scaling.
    \item \textbf{Database Performance Considerations:} The system shall include configuration guidelines and optional node affinity/IO provisioning primitives (e.g., hostPath, CSI volumes with provisioned IOPS) to mitigate containerized I/O bottlenecks for production database workloads.
\end{itemize}

\subsubsection*{Flexibility (Cloud-Agnostic Design)}
\begin{itemize}
    \item \textbf{Vendor Lock-In Mitigation:} The system architecture shall be cloud-agnostic, supporting deployment across multiple providers (e.g., AWS, GCP, Azure).
    \item \textbf{Portability:} By relying on standards-based technologies such as containers and Kubernetes instead of proprietary cloud APIs, the system shall allow seamless workload migration with minimal refactoring.
    \item \textbf{Strategic Value:} This approach mitigates hyperscaler walled-garden strategies that rely on proprietary services and data egress fees.
\end{itemize}

\subsubsection*{Security \& Monitoring}
\begin{itemize}
    \item \textbf{Access Control:} The system shall implement robust authentication and authorization mechanisms (e.g., OAuth/OIDC, RBAC) to manage user access to database resources.
    \item \textbf{Tenant Isolation:} Strong tenant isolation shall be enforced to prevent unauthorized access to sensitive data in multi-tenant environments (network policies, namespace separation, encrypted volumes).
    \item \textbf{Administrative Monitoring:} An administrative dashboard shall be provided to monitor system performance, user activity, tool/agent usage, and resource utilization to support compliance and proactive management.
    \item \textbf{Agent Safety Controls:} The platform shall enforce agent capability scoping (which agents can execute which tools), rate limits on agent-initiated operations, and verification layers (schema validators, transaction dry-runs) before executing critical changes.
    \item \textbf{Auditability:} All agent interactions, approvals, tool calls, and schema/data mutations shall be logged with cryptographic timestamps to support forensic analysis and regulatory audits.
\end{itemize}
